\documentclass[10pt,a4paper]{amsart}
\usepackage[margin=19mm]{geometry}
\usepackage{amsmath,amssymb,mathtools}
\usepackage[scr=rsfs,cal=euler]{mathalfa}
\usepackage{xfrac}
\usepackage[unicode,hidelinks,bookmarks=false]{hyperref}
\newcommand{\ie}{i.e.\ }
\newcommand{\cf}{cf.\ }
\newcommand{\resp}{respectively\ }
\newcommand{\Subsection}{\subsection}
\providecommand{\nobreakdash}{\nobreak\hbox{-}}
\setlength{\emergencystretch}{2em}
\multlinegap0pt
\usepackage{bm}
\usepackage{bbm}
\usepackage{dsfont}
\usepackage{xspace}
\usepackage{cancel}
\usepackage{mathtools}
\usepackage{amsthm}
\makeatletter
\@ifundefined{lemma}{\newtheorem{lemma}{Lemma}}{}
\makeatother
\title[Fronthaul Compression for Uplink Cloud-RAN with Finite-Alpha]{Fronthaul Compression for Uplink Cloud-RAN with Finite-Alphabet Inputs: A Reverse Mercury/Waterfilling Approach}
\author{Subin Shin, Jaehoon Lee, Seok-Hwan Park,, Jeonghun Park}
\date{Source: arXiv:2609.21265v1 (CC BY 4.0)}
\begin{document}
\maketitle

\noindent\emph{Source and licence.} Fronthaul Compression for Uplink Cloud-RAN with Finite-Alphabet Inputs: A Reverse Mercury/Waterfilling Approach. Version arXiv:2609.21265v1 (CC BY 4.0), distributed under the Creative Commons Attribution 4.0 International licence, as recorded in the arXiv licence metadata for this exact version and shown on the abstract page https://arxiv.org/abs/2609.21265v1. Full rights notice: licenses/arxiv-2609.21265-CC-BY-4.0.txt.\par\smallskip

\noindent\textit{Editorial scope.} Counted unit: the Lemma whose source label is \texttt{lem:dsinr} in the article (source0.tex lines 1221--1237, source label \texttt{lem:dsinr}), delivered together with the article's own complete original proof of it (source0.tex lines 1239--1264, one \texttt{proof} environment of 26 lines). No mathematics is written, repaired or re-derived: statement and proof are the article's own text line for line. Required context retained uncounted: none. Every definition, display and label of those lines is retained unchanged. Omitted from this package and not redistributed: 1--1220, 1238--1238, 1265--2110. Nothing inside a retained excerpt is omitted or altered: every retained body is delivered exactly as the article's lines are written, so no omission receipt of any kind accompanies this package. Citations: the retained excerpts carry no citation command at all, so no bibliography entry of the article is transcribed and no reference is redistributed. Reference adaptation: none was needed and none was made; every reference inside the delivered unit resolves inside it, so no unresolved reference or citation is produced. Printed slips kept literal: no printed word, symbol, hypothesis or number of the counted statement or of its proof was altered. Layout and class: the article's own class is \texttt{IEEEtran} with packages including fontenc, graphicx, epsfig, cite, mcite, amsfonts, helvet, fancyhdr, threeparttable, epsf, amsthm, amsmath, siunitx, amssymb; the wrapper is the lane's pinned \texttt{amsart} editorial wrapper at 10pt on A4, which restates the theorem-like environments the retained text needs and adds the article's own macro definitions that the retained text needs (none), with extra wrapper packages (bm, bbm, dsfont, xspace, cancel, mathtools). The delivered numbering is the unit's own. Assets: no retained span contains a figure environment, a graphics-inclusion command, a PDF-page inclusion, a table or a TikZ picture; no image file is copied into this package, the package declares no asset dependency and no image file is redistributed with it. Licence: version arXiv:2609.21265v1 of this article is distributed under the Creative Commons Attribution 4.0 International licence, as recorded in the arXiv licence metadata for this exact version and shown on the abstract page https://arxiv.org/abs/2609.21265v1. A full rights notice is delivered at licenses/arxiv-2609.21265-CC-BY-4.0.txt. Characters: every retained excerpt is pure ASCII, so the delivered engine, XeLaTeX with its default Latin Modern text font, reports no missing character.
\begin{lemma}[Derivative of post-LMMSE SINR]
\label{lem:dsinr}
Let $\mathrm{SINR}_k = P_k \mathbf{h}_k^{\sf H} 
\mathbf{R}_k^{-1} \mathbf{h}_k$, where $\mathbf{R}_k \triangleq \sum_{j \neq k} P_j \mathbf{h}_j \mathbf{h}_j^{\sf H} + \sigma^2 \mathbf{I}_{MN} + \mathbf{Q}$. 
Then,
\begin{align} \label{eq:dsinr}
  \frac{\partial \mathrm{SINR}_k}{\partial b_{m,n}} 
  = (\ln 2) \, q_{m,n} \cdot P_k \, 
    \big|\mathbf{h}_k^{\sf H} \mathbf{R}_k^{-1} 
         \mathbf{\tilde u}_{m,n}\big|^2,
\end{align}
% where $\mathbf{u}_{m,n} \in \mathbb{C}^{MN}$ 
% denotes the $n$-th eigenvector of RU $m$ embedded 
% in the aggregated coordinate system (i.e., zero 
% outside the $m$-th block).
where $\mathbf{\tilde u}_{m,n} \triangleq \mathbf{v}_m \otimes \mathbf{u}_{m,n} \in \mathbb{C}^{MN}$, with $\mathbf{u}_{m,n} \in \mathbb{C}^N$ being the $n$-th eigenvector of the local covariance matrix $\mathbf{R}_{\mathbf{y}_m}$ at RU $m$, and $\mathbf{v}_m \in \mathbb{R}^M$ denoting the $m$-th standard basis vector.
\end{lemma}

\begin{proof}
The dependence of $\mathrm{SINR}_k$ on $b_{m,n}$ is 
mediated entirely through the quantization noise 
covariance $\mathbf{Q}$. By the chain rule,
\begin{align}
  \frac{\partial \mathrm{SINR}_k}{\partial b_{m,n}} 
  = P_k \mathbf{h}_k^{\sf H} 
    \frac{\partial \mathbf{R}_k^{-1}}
         {\partial b_{m,n}} \mathbf{h}_k.
\end{align}
The matrix derivative identity 
$\partial \mathbf{R}_k^{-1}/\partial b_{m,n} 
= -\mathbf{R}_k^{-1} (\partial \mathbf{R}_k/
\partial b_{m,n}) \mathbf{R}_k^{-1}$ together with 
$\partial \mathbf{R}_k/\partial b_{m,n} 
= \partial \mathbf{Q}/\partial b_{m,n} 
= -(\ln 2) q_{m,n} \mathbf{\tilde u}_{m,n} \mathbf{\tilde u}_{m,n}^{\sf H}$, where the last equality follows from $q_{m,n} = c \lambda_{m,n} \cdot 2^{-b_{m,n}}$ and the block structure of $\mathbf{Q}$ gives
\begin{align}
  \frac{\partial \mathbf{R}_k^{-1}}{\partial b_{m,n}} 
  = (\ln 2) q_{m,n} \mathbf{R}_k^{-1} 
    \mathbf{\tilde u}_{m,n} \mathbf{\tilde u}_{m,n}^{\sf H} 
    \mathbf{R}_k^{-1}.
\end{align}
Substituting this into the chain rule expression 
yields~\eqref{eq:dsinr}.
\end{proof}

\end{document}
